\documentclass[12pt]{article}

\usepackage{graphicx}
\usepackage{float}
\usepackage{booktabs}
\usepackage{listings}
\usepackage{xcolor}
\usepackage{geometry}

\geometry{margin=1in}

\title{Distributed and Parallel Computing \\ Lab 3 Report \\ MPI Collective Operations}
\author{Ahmad Bin Rashid}
\date{}

\lstset{
language=C,
basicstyle=\ttfamily\small,
keywordstyle=\color{blue},
commentstyle=\color{gray},
stringstyle=\color{red},
numbers=left,
numberstyle=\tiny,
frame=single,
breaklines=true
}

\begin{document}

\begin{titlepage}

\title{
    \vspace{0.5cm}
    \textbf{\Large Lab 3 Report} \\
    \vspace{0.5cm}
    \large MPI Collective Operations \\
    \vspace{1.5cm}
}

\author
{\Large \textbf{Name:} Ahmad Bin Rashid \\ [0.5cm]\\
\Large \textbf{Roll No:} 2023-CS-53 \\ [0.5cm]\\
\Large \textbf{Instructor:} Sir Waqas Ali \\ [0.5cm]\\
\Large \textbf{Course:} Distributed and Parallel Computing \\ [0.5cm]\\ 
\Large \textbf{University of Engineering and Technology} \\ [0.5cm]\\ }

\begin{figure}
    \centering
    \vspace{1cm}
    \includegraphics[width=0.3\textwidth]{uet-lahore-seeklogo.png}
\end{figure}

\date{\today}

\end{titlepage}

% --- Title Page ---
\maketitle
\thispagestyle{empty}
\newpage

\section{Exercise Execution Screenshots}

\subsection{Exercise 1: MPI Broadcast}

\begin{figure}[H]
\centering
\fbox{\includegraphics[width=0.9\textwidth]{Screenshot 2026-03-12 115439.png}}
\caption{Successful execution of the MPI Broadcast example}
\end{figure}

\subsection{Exercise 2: MPI Reduce}

\begin{figure}[H]
\centering
\fbox{\includegraphics[width=0.9\textwidth]{Screenshot 2026-03-12 115449.png}}
\caption{Execution output for MPI Reduce example}
\end{figure}

\subsection{Exercise 3: MPI Scatter and Gather}

\begin{figure}[H]
\centering
\fbox{\includegraphics[width=0.9\textwidth]{Screenshot 2026-03-12 115458.png}}
\caption{Execution output for Scatter and Gather example}
\end{figure}

\subsection{Exercise 4: MPI Scatterv and Gatherv}

\begin{figure}[H]
\centering
\fbox{\includegraphics[width=0.9\textwidth]{Screenshot 2026-03-12 115515.png}}
\caption{Execution output demonstrating uneven distribution using Scatterv/Gatherv}
\end{figure}

\subsection{Exercise 5: Matrix Vector Multiplication}

\begin{figure}[H]
\centering
\fbox{\includegraphics[width=0.9\textwidth]{Screenshot 2026-03-12 115525.png}}
\caption{Matrix-vector multiplication result using MPI}
\end{figure}


\section{Code Listing for Exercise 4 Modifications}

The following code modifications were made to correctly implement
\texttt{MPI\_Scatterv} and \texttt{MPI\_Gatherv} for uneven data distribution.

\begin{lstlisting}
// Key modifications for Scatterv/Gatherv implementation

// Fix broken declarations
int *sendcounts;
int *displs;

// Compute sendcounts and displacements for all processes
for (int i = 0; i < size; i++) {
    sendcounts[i] = N / size;
    if (i < N % size)
        sendcounts[i]++;

    displs[i] = sum;
    sum += sendcounts[i];
}

// Scatter uneven chunks
MPI_Scatterv(data, sendcounts, displs, MPI_INT,
             local_data, local_n, MPI_INT,
             0, MPI_COMM_WORLD);

// Gather results
MPI_Gatherv(local_data, local_n, MPI_INT,
            gathered_data, sendcounts, displs,
            MPI_INT, 0, MPI_COMM_WORLD);
\end{lstlisting}

\textbf{Explanation of Fixes and Additions:}

\begin{itemize}
\item Fixed broken variable declarations caused by comment formatting errors.
\item Corrected the loop responsible for calculating \texttt{sendcounts} and \texttt{displs}.
\item Ensured that all processes compute their own \texttt{local\_n} to avoid an extra \texttt{MPI\_Scatter}.
\item Added \texttt{MPI\_Gatherv} to collect uneven chunks back to the root process.
\item Implemented verification to confirm gathered data matches the original dataset.
\end{itemize}


\section{Performance Comparison for MPI Reduce}

\begin{table}[H]
\centering
\begin{tabular}{ccc}
\toprule
Processes & Execution Time (seconds) & Result \\
\midrule
1 & 0.000005 & 500500 \\
2 & 0.000011 & 500500 \\
4 & 0.000029 & 500500 \\
8 & 0.000046 & 500500 \\
\bottomrule
\end{tabular}
\caption{Execution time comparison for MPI Reduce example}
\end{table}


\section{Discussion: Collective vs Point-to-Point Communication}

Collective operations in MPI provide a higher-level abstraction for common
communication patterns such as broadcasting, reduction, and data distribution.
Unlike point-to-point communication, where the programmer must explicitly
manage message exchanges between processes, collective operations allow
multiple processes to participate in a coordinated communication step.

These operations are usually optimized internally by the MPI implementation,
resulting in better scalability and performance. They also simplify program
development by reducing the complexity of synchronization and message
management. In contrast, point-to-point communication offers more flexibility
but requires careful coordination to avoid issues such as deadlocks and
inefficient communication patterns.

\section{Self-Assessment Answers}

\textbf{1. What is the main difference between point-to-point and collective communication?} \\
Point-to-point communication involves explicit message exchange between two specific processes using functions like \texttt{MPI\_Send} and \texttt{MPI\_Recv}. Collective communication involves a group of processes participating in a coordinated operation such as \texttt{MPI\_Bcast} or \texttt{MPI\_Reduce}.

\vspace{0.3cm}

\textbf{2. Why must all processes call a collective operation like MPI\_Bcast?} \\
All processes in the communicator must call the collective function because MPI synchronizes them internally. If some processes skip the call, the program may hang or produce undefined behavior.

\vspace{0.3cm}

\textbf{3. In the matrix-vector multiplication example, why do we broadcast \texttt{x} instead of scattering it?} \\
The vector \texttt{x} is needed entirely by every process to compute its portion of the result vector. Broadcasting ensures all processes receive the full vector.

\vspace{0.3cm}

\textbf{4. Write the MPI\_Reduce call that computes the maximum value across all processes.} \\
\begin{lstlisting}
MPI_Reduce(&local_value, &max_value, 1, MPI_INT, MPI_MAX, 0, MPI_COMM_WORLD);
\end{lstlisting}

\vspace{0.3cm}

\textbf{5. Which collective would you use to collect processed results back to process 0?} \\
\texttt{MPI\_Gather} would be used to collect the processed data from all processes and store it in a receive buffer on process 0.

\vspace{0.3cm}

\textbf{6. What happens if the root in MPI\_Gather provides a receive buffer that is too small?} \\
The program may cause memory corruption, incorrect results, or a runtime error because incoming data exceeds the allocated buffer.

\vspace{0.3cm}

\textbf{7. In the scatterv example, explain the purpose of \texttt{sendcounts} and \texttt{displs}.} \\
\texttt{sendcounts} specifies how many elements each process receives, while \texttt{displs} specifies the starting position of each chunk in the original array.

\vspace{0.3cm}

\textbf{8. Why is MPI\_Reduce usually faster than manual MPI\_Send/MPI\_Recv implementations?} \\
MPI collective operations use optimized algorithms (such as tree-based reductions) that minimize communication steps, making them more efficient than manual implementations.

\vspace{0.3cm}

\textbf{9. How would you modify the matrix-vector multiplication to use MPI\_Allgather?} \\
Each process would compute its partial result and use \texttt{MPI\_Allgather} instead of \texttt{MPI\_Gather} so that every process receives the complete result vector \texttt{y}.

\end{document}